\scanpage{84}
i.e. on a
\[ \sigma_0=\begin{pmatrix}-1&\deleted{\unsure{1+}}\alpha&1\\0&\deleted{\unsure{2}}1+\beta&0\\0&0&1\end{pmatrix} \tag*{(17)} \]
\[ \sigma_0(x,y)=(1-x+\deleted{\unsure{(1+}}\alpha y,(1+\beta)y) \tag*{(18)} \]
\[ \left\{\begin{aligned}&\left\{\begin{aligned}\sigma_1s_0&=s_0\\
\sigma_1s_1&=s_{-1}\\
\sigma_1s_{-1}&=s_1\end{aligned}\right.\\
&\left\{\begin{aligned}\sigma_1e_1&=e_2\\
\sigma_1e_2&=e_1\end{aligned}\right.\end{aligned}\right.\qquad\text{i.e.} \tag*{(\deleted{\unsure{18}}19)} \]
\[ \sigma_1=\begin{pmatrix}0&1&0\\1&0&0\\0&0&1\end{pmatrix} \tag*{(20)} \]
\[ \sigma_1(x,y)\unsure{{:}|{=}}(y,x) \tag*{(21)} \]

Il reste à écrire la condition $\sigma_0^2=\operatorname{id}$.
On trouve
\[ \left\{\begin{aligned}&\boxed{\alpha\beta=0}\qquad\deleted{(1+\beta)^2=1}\\
&(1+\beta)^2=1\quad\text{i.e}\quad\boxed{\beta^2+2\beta=0}\quad\text{i.e}\quad\boxed{\beta(\beta+2)=0}\end{aligned}\right. \tag*{\deleted{\unsure{(22)}}} \]
Donc :

\begin{tabular}{|p{.91\linewidth}}\uline{Proposition} Les représentations géom affines
régulières de $\Pi_\infty$ \inserted{sur $k$} sont \deleted{définies} en corr.
1-1 \unsure{can.|corr.} avec les couples $(\alpha,\beta)$ d'éléments
de $k$ satisfaisant les conditions
\[ \alpha\beta=(\beta+2)\beta=0 \tag*{(22)} \]
\end{tabular}
% archive sheet number: 83
